% Part: first-order-logic
% Chapter: models-theories
% Section: theories

\documentclass[../../../include/open-logic-section]{subfiles}

\begin{document}

\olfileid[pt]{fol}{mat}{the}

\olsection{Exemplos de Teorias de Primeira Ordem}

\begin{ex}
A teoria das ordens lineares estritas na linguagem~$\Lang L_<$ é
axiomatizada pelo conjunto
\begin{align*}
\{\quad & \lforall[x][\lnot x < x], \\
& \lforall[x][\lforall[y][((x < y \lor y <
    x) \lor x = y)]], \\
& \lforall[x][\lforall[y][\lforall[z][((x < y
      \land y < z) \lif x < z)]]] \quad \}
\end{align*}
Ela captura completamente as !!{structure}s pretendidas: toda ordem linear
estrita é um modelo deste sistema de axiomas e, vice-versa, se $R$
é uma ordem linear sobre um conjunto $X$, então a estrutura $\Struct M$ com
$\Domain M = X$ e $\Assign{<}{M} = R$ é um modelo desta teoria.
\end{ex}

\begin{ex}
A teoria dos grupos na linguagem $\Obj 1$ (!!{constant}), $\cdot$
(!!{function} de dois lugares) é axiomatizada por
\begin{align*}
& \lforall[x][\eq[(x \cdot \Obj 1)][x]]\\
& \lforall[x][\lforall[y][\lforall[z][\eq[(x \cdot (y \cdot z))][((x
          \cdot y) \cdot z)]]]]\\
& \lforall[x][\lexists[y][\eq[(x \cdot y)][\Obj 1]]]
\end{align*}
\end{ex}

\begin{ex}
A teoria da aritmética de Peano é axiomatizada pelas seguintes
sentenças na linguagem da aritmética~$\Lang L_A$.
\begin{align*}
& \lforall[x][\lforall[y][(\eq[x'][y'] \lif \eq[x][y])]]\\
& \lforall[x][\eq/[\Obj 0][x']]\\
& \lforall[x][\eq[(x + \Obj 0)][x]]\\
& \lforall[x][\lforall[y][\eq[(x + y')][(x + y)']]]\\
& \lforall[x][\eq[(x \times \Obj 0)][\Obj 0]]\\
& \lforall[x][\lforall[y][\eq[(x \times y')][((x \times y) + x)]]]\\
& \lforall[x][\lforall[y][(x < y \liff \lexists[z][\eq[(z' + x)][y])]]]\\
\intertext{mais todas as sentenças da forma}
& (!A(\Obj 0) \land \lforall[x][(!A(x) \lif !A(x'))]) \lif \lforall[x][!A(x)]
\end{align*}
Como há infinitas sentenças desta última forma, este
sistema de axiomas é infinito.  Esta última forma é chamada de
\emph{esquema de indução}. (Na verdade, o esquema de indução é um pouco mais
complicado do que deixamos transparecer aqui.)

O último axioma é uma \emph{definição explícita} de~$<$.
\end{ex}

\begin{ex}
A teoria dos conjuntos puros desempenha um papel importante nos fundamentos
(e na filosofia) da matemática.  Um conjunto é puro se todos os seus
!!{element}s são também conjuntos puros.  O conjunto vazio conta, portanto, como
puro, mas um conjunto que tem como !!a{element} algo que não é um conjunto
não seria puro.  Então os conjuntos puros são aqueles que são formados apenas
a partir do conjunto vazio e sem ``urelementos'', isto é, objetos que não são,
eles próprios, conjuntos.

O seguinte poderia ser considerado um sistema de axiomas para uma teoria dos
conjuntos puros:
\begin{align*}
& \lexists[x][\lnot \lexists[y][y \in x]]\\
& \lforall[x][\lforall[y][(\lforall[z](z \in x \liff z \in y) \lif 
\eq[x][y])]]\\
& \lforall[x][\lforall[y][\lexists[z][\lforall[u][(u \in z \liff
(\eq[u][x] \lor \eq[u][y]))]]]]\\
& \lforall[x][\lexists[y][\lforall[z][(z \in y \liff \lexists[u][(z \in
        u \land u \in x)])]]]\\
\intertext{mais todas as sentenças da forma} &
\lexists[x][\lforall[y][(y \in x \liff !A(y))]]
\end{align*}
O primeiro axioma diz que há um conjunto sem !!{element}s (isto é,
$\emptyset$ existe); o segundo diz que conjuntos são extensionais; o
terceiro que, para quaisquer conjuntos $X$ e $Y$, o conjunto $\{X, Y\}$ existe; o
quarto que, para qualquer conjunto $X$, o conjunto $\cup X$ existe, onde $\cup X$ é a 
união de todos os elementos de $X$.

As !!{sentence}s mencionadas por último são chamadas coletivamente de
\emph{esquema de compreensão ingênuo}.  Ele essencialmente diz que, para todo
$!A(x)$, o conjunto $\Setabs{x}{!A(x)}$ existe---então, à primeira vista, um
axioma verdadeiro, útil e talvez até necessário.  É chamado de ``ingênuo''
porque, como se vê, ele torna esta teoria insatisfatível: se você
toma $!A(y)$ como $\lnot y \in y$, obtém a !!{sentence}
\[
\lexists[x][\lforall[y][(y \in x \liff \lnot y \in y)]]
\]
e esta !!{sentence} não é satisfeita em nenhuma !!{structure}.
\end{ex}

\begin{ex}
Na área da \emph{mereologia}, a relação de \emph{parte} é uma
relação fundamental.  Assim como as teorias de conjuntos, há teorias
de parte que axiomatizam várias concepções (às vezes
conflitantes) dessa relação.

A linguagem da mereologia contém um único símbolo de predicado
de dois lugares~$\Obj P$, e $\Atom{\Obj P}{x, y}$ ``significa'' que $x$ é uma
parte de~$y$.  Quando temos essa interpretação em mente, !!a{structure}
para essa linguagem é chamada de \emph{estrutura de parte}.  É claro que
nem toda estrutura para um único predicado de dois lugares realmente
merecerá esse nome.  Para ter uma chance de capturar ``parte'',
$\Assign{\Obj P}{M}$ deve satisfazer algumas condições, que podemos
estabelecer como axiomas para uma teoria de parte.  Por exemplo, a parte é uma
ordem parcial sobre objetos: todo objeto é uma parte (ainda que uma
parte \emph{imprópria}) de si mesmo; dois objetos diferentes não podem ser partes
um do outro; uma parte de uma parte de um objeto é, ela própria, parte daquele
objeto.  Note que, nesse sentido, ``é uma parte de'' se assemelha a ``é um
subconjunto de'', mas não se assemelha a ``é um elemento de'', que não é
nem reflexiva nem transitiva.
\begin{align*}
& \lforall[x][\Atom{\Obj P}{x,x}] \\
& \lforall[x][\lforall[y][((\Part{x}{y} \land \Part{y}{x})
      \lif \eq[x][y])]] \\
& \lforall[x][\lforall[y][\lforall[z][((\Part{x}{y} \land
        \Part{y}{z}) \lif \Part{x}{z})]]]\\
\intertext{Além disso, quaisquer dois objetos têm uma soma mereológica (um objeto que tem
  esses dois objetos como partes, e é minimal nesse aspecto).}  &
\lforall[x][\lforall[y][\lexists[z][\lforall[u][(\Part{z}{u} \liff
        (\Part{x}{u} \land \Part{y}{u}))]]]]
\end{align*}
Estes são apenas alguns dos princípios básicos de parte considerados pelos
metafísicos.  Princípios adicionais, no entanto, rapidamente se tornam difíceis
de formular ou escrever sem primeiro introduzir algumas relações
definidas.  Por exemplo, a maioria dos metafísicos interessados em mereologia
também considera o seguinte como um princípio válido: sempre que um
objeto~$x$ tem uma parte própria~$y$, ele também tem uma parte~$z$ que não tem
partes em comum com~$y$, e tal que a fusão de $y$ e $z$ é
$x$.
\end{ex}

\end{document}
